\DocumentMetadata{
  pdfversion=2.0,
  pdfstandard=ua-2,
  lang=en-US,
  tagging=on,
  tagging-setup={math/setup=mathml-SE,tabsorder=structure}
}
\documentclass[11pt,letterpaper]{article}
\usepackage[margin=0.68in,includefoot]{geometry}
\usepackage{newcomputermodern}
\usepackage{amsmath,amscd,mathtools}
\usepackage{tikz,adjustbox}
\providecommand{\square}{\mdlgwhtsquare}
\usepackage[hidelinks]{hyperref}
\hypersetup{
  pdftitle={Combinatorics PhD Exam, August 2019},
  pdfauthor={University of Florida Department of Mathematics},
  pdfsubject={Graduate mathematics examination},
  pdfdisplaydoctitle=true,
  bookmarksopen=true
}
\setcounter{secnumdepth}{0}
\setlength{\parindent}{0pt}
\setlength{\parskip}{0.28em}
\setlength{\emergencystretch}{3em}
\setlength{\leftmargini}{2.15em}
\setlength{\leftmarginii}{2.65em}
\setlength{\leftmarginiii}{3em}
\renewcommand{\labelenumi}{\textbf{\theenumi.}}
\pagestyle{empty}
\raggedbottom
\allowdisplaybreaks
\newenvironment{examproblems}[1]{%
  \begin{enumerate}%
  \setcounter{enumi}{#1}%
  \setlength{\topsep}{0.35em}%
  \setlength{\partopsep}{0pt}%
  \setlength{\itemsep}{0.48em}%
  \setlength{\parsep}{0.18em}%
}{\end{enumerate}}
\newenvironment{examsubparts}{%
  \begin{enumerate}%
  \setlength{\topsep}{0.18em}%
  \setlength{\partopsep}{0pt}%
  \setlength{\itemsep}{0.16em}%
  \setlength{\parsep}{0.1em}%
}{\end{enumerate}}
\newenvironment{examinstructions}{%
  \par\begingroup\itshape
}{%
  \par\endgroup\vspace{0.18em}
}
\makeatletter
\renewcommand\section{\@startsection{section}{1}{0pt}%
  {-1.25ex plus -.25ex minus -.1ex}%
  {0.55ex plus .1ex}%
  {\normalfont\large\bfseries}}
\renewcommand{\@maketitle}{%
  \newpage\null\vskip -1em%
  {\footnotesize\bfseries
    \href{https://www.ufl.edu/}{University of Florida}, \href{https://math.ufl.edu/}{Department of Mathematics}\par}%
  \vskip -0.5em%
  \tagstructbegin{tag=Title}%
    {\LARGE\bfseries\href{https://gma.math.ufl.edu/exams/combinatorics-phd/}{Combinatorics PhD Exam}, August 2019\par}%
  \tagstructend%
  \vskip 0.25em%
  \tagmcbegin{artifact}%
    \hrule height 1pt%
  \tagmcend%
  \vskip 1em%
}
\makeatother
\title{Combinatorics PhD Exam, August 2019}
\author{}
\date{}
\begin{document}
\maketitle
\thispagestyle{empty}
\begin{examproblems}{0}
\item
List the infinite families of symmetries of the standard square tiling of the plane. (No two families should have more than the identity in common.)
\item
Let \((v,b,r,k,\lambda)\) denote the number, respectively, of points, blocks, blocks containing a given point, points in each block, and blocks containing two distinct points of a balanced incomplete block design. Prove that 
\[
bk=vr
\]
 and 
\[
\lambda(v-1)=r(k-1).
\]
\item
Let \(h_n\) be the number of all words over the alphabet \(\{A,B,C\}\) that are of length~\(n\) and in which a letter~\(A\) is never immediately followed by a letter~\(B\). Set \(h_0=1\). Find a closed form for the ordinary generating function 
\[
H(z)=\sum_{n\ge 0}h_nz^n.
\]
\item
Let \(\ell_n\) be the total number of leaves in all unlabeled plane binary trees on~\(n\) vertices. (These trees are rooted, and each vertex has at most two children, and every child is either a left child or a right child, even if it is an only child.) Note that \(\ell_1=1\), \(\ell_2=2\), and \(\ell_3=6\).
\begin{examsubparts}
\item[\textnormal{(a)}]
Find the generating function 
\[
L(z)=\sum_{n\ge 1}\ell_nz^n
\]
 in closed form.
\item[\textnormal{(b)}]
Find an explicit formula for the numbers \(\ell_n\).
\end{examsubparts}
\item
Let \(t_n\) be the number of ways to choose a permutation of length~\(n\), and then to color a subset (not necessarily proper) of its even cycles red. Find a closed formula for the exponential generating function of the numbers \(t_n\).
\item
Recall that for a prime power~\(q\), the Gaussian or~\(q\)-binomial coefficient \(\left[\begin{smallmatrix}n\\k\end{smallmatrix}\right]_q\) is defined to be the number of~\(k\) dimensional subspaces of an~\(n\)-dimensional vector space over \(\mathrm{GF}(q)\). Equivalently, we can think of \(\left[\begin{smallmatrix}n\\k\end{smallmatrix}\right]_q\) as counting \(k\times n\) matrices over \(\mathrm{GF}(q)\) which are in reduced row echelon form and have no all-zero rows. Using either of these interpretations, prove that 
\[
\left[\begin{smallmatrix}n\\k\end{smallmatrix}\right]_q
=
\left[\begin{smallmatrix}n-1\\k\end{smallmatrix}\right]_q
+q^{n-k}\left[\begin{smallmatrix}n-1\\k-1\end{smallmatrix}\right]_q
\]
 for all \(n,k\ge 1\).
\item
Recall that \(R(k,\ell)\) is the minimum integer~\(n\) such that every graph on~\(n\) vertices contains either a clique (complete subgraph) on~\(k\) vertices or an independent set on~\(\ell\) vertices. Prove that \(R(k,\ell)\) is finite for all~\(k\) and~\(\ell\).
\item
Prove that \(R(k,k)>2^{k/2}\) for all~\(k\). Hint: Consider a random graph.
\end{examproblems}
\end{document}
